\section{Conclusion}
\label{sec:conclusion}

We have presented an empirical study of self-training unsupervised
domain adaptation for cross-domain Mongolian text detection.
On a clean two-domain bilingual benchmark we diagnose three failure
modes that the prevailing UDA recipe overlooks when scaled down to
small corpora and constrained compute. First, we identify a
\emph{cold-start collapse} in which a randomly initialised teacher
produces noise pseudo-labels that the student cannot escape, and show
that a one-line remedy -- initialising both networks from the
source-only checkpoint -- raises the harder of our two transfer
directions from zero true positives to a measurable
$\hmean = 0.0654$ after threshold sweeping. Second, we document an
\emph{asymmetric UDA gain phenomenon}: the same pipeline yields a
$+6.26$~pp $\hmean$ jump where source-only is structurally broken,
and only $+0.14$~pp where source-only already produces a weak target
signal. Third, we show that this pipeline has a structural
\emph{5-epoch sweet spot}; extending training to 25 epochs collapses
$\hmean$ in both directions through teacher--student drift, even
under a conservative pseudo-label weight schedule.

The headline message of the work is that for under-resourced
cross-domain text detection on a constrained compute budget,
self-training UDA is most useful precisely where the source-only
baseline is most broken, and stays useful only within a narrow
training window. We hope this regime-aware perspective will guide
both method design and reporting standards in future low-budget UDA
work.
